%==========================================================================
%  TRES0312 Fysikk — Foiler-mal
%
%  MØrk/lys modus: kommenter ut én av de to linjene nedenfor.
%==========================================================================
\newif\ifdarkmode
\darkmodetrue      % ← mørk modus
%\darkmodefalse   % ← lys modus  (fjern % foran og legg % foran linjen over)

\documentclass[aspectratio=169, 12pt]{beamer}

\usepackage{amd-beamer}

%--------------------------------------------------------------------------
%  DOKUMENTINFO
%--------------------------------------------------------------------------
\title{Velkommen til fysikk!}
\subtitle{TRES0312 --- Introduksjon}
\author{Ben David Normann}
\date{\today}

%--------------------------------------------------------------------------
%  SEKSJONSKONTEKST
%  \setcounter{section}{N} setter gjeldende seksjonsnummer til N, slik at
%  \defn, \exm og \oppgave får samme nummerering som i kompendiet.
%  Tellerne nullstilles automatisk ved hvert \setcounter{section}{...}.
%  Bruk gjerne \setcounter{section}{N} rett før en frame der du bytter
%  seksjon, i stedet for én gang globalt.
%--------------------------------------------------------------------------
\setcounter{section}{1}   % Seksjon 1: "Hva er fysikk?"

\begin{document}

%--------------------------------------------------------------------------
%  SLIDE 1 — TITTELSIDE
%--------------------------------------------------------------------------
\begin{frame}[plain]
  \titlepage
\end{frame}

%--------------------------------------------------------------------------
%  SLIDE 2 — TOM
%--------------------------------------------------------------------------
\begin{frame}{}
  
  \spq{Hva er fysikk?}
  \pause
  \svar{
  Fysikk er studiet av naturens grunnstruktur, gjerne med matematikk som verktøy.
  }
\pause
  \tanke{}{%
    Kan alt i verden beskrives i matematikk? Hva med \textit{kjærlighet} eller \textit{bevissthet}? Slike vanskelige, filosofiske spørsmål skal vi ikke vsare på i dette kompendiet, men de er verdt å merke seg.}
 
\end{frame}

%--------------------------------------------------------------------------
%  SLIDE 3 — FYSISKE STØRRELSER
%--------------------------------------------------------------------------

\begin{frame}{Fysiske størrelser}
  \spq{Hvordan beskriver vi virkeligheten ved hjelp av matematikk?}
  \pause
  \svar{
    \begin{itemize}
      \item Vi bruker fysiske størrelser.
      \item Vi må være nøyaktige
    \end{itemize}
  }
  \pause
  \defnxr{fys-storrelse}{Fysisk størrelse}{%
    En \textit{fysisk størrelse} har både \textit{verdi} og \textit{enhet}:
    \[
      \textrm{Fysisk størrelse} = \textrm{verdi} \cdot \textrm{enhet}.
    \]
  }

\end{frame}

%--------------------------------------------------------------------------
%  SLIDE 4 — Enheter
%--------------------------------------------------------------------------

\begin{frame}{Enheter}
  \vspace{3em}
\spq{Hvor mange forskjellige enheter trenger man for å beskrive \textit{alt} i naturen?}
  \pause
  \svar{7}

  

\end{frame}

%--------------------------------------------------------------------------
%  SLIDE 5 — SI-ENHETENE i sammenheng 
%--------------------------------------------------------------------------

\begin{frame}{SI-enhetene}
  \begin{center}
    \includegraphics[width=0.56\textwidth, height=0.72\textheight, keepaspectratio]{Kapittel1:introduksjon/SI-enheter_transp.png}\\[0.2em]
    {\small SI-systemet}
  \end{center}
\end{frame}


%--------------------------------------------------------------------------
%  SLIDE 6 — SI-ENHETENES NATURKONSTANTER
%--------------------------------------------------------------------------

\begin{frame}{SI-enhetene og naturkonstantene}
  \begin{center}
    \includegraphics[width=0.56\textwidth, height=0.7\textheight, keepaspectratio]{Kapittel1:introduksjon/SI-enheter2_transp.png}\\[0.2em]
    {\small Den ytre ringen viser naturkonstantene som definerer de sju SI-enhetene (2018-definisjonen).}
  \end{center}
\end{frame}

%--------------------------------------------------------------------------
%  SLIDE 7 — Noen enheter...
%--------------------------------------------------------------------------

\begin{frame}{}
\renewcommand{\arraystretch}{1.3}
\aktivitet{Hva er enhetene for disse fysiske størrelsene?}{%
\begin{center}
\begin{tabular}{|l|c|l|}
\hline
\rowcolor{blue!10!black!80}
\textbf{Fysisk størrelse} & \textbf{Enhet} & \textbf{Beskrivelse} \\
\hline
Tid &  & Hvor lenge noe varer \\
Lengde &  & Avstand mellom to punkter \\
Masse &  & Hvor mye stoff det er i et objekt \\
Fart &  & Hvor raskt noe beveger seg \\
Akselerasjon &  & Hvor raskt farten endrer seg \\
Kraft &  & Vekselvirkning mellom legemer \\
Energi &  & Potensiale for vekselvirkning. \\
\hline
\end{tabular}
\end{center}
}
\end{frame}


%--------------------------------------------------------------------------
%  SLIDE 8 — Oppgave 1.1
%--------------------------------------------------------------------------

\begin{frame}{Oppgave}
  \oppgavexr{si-enheter}{SI-enheter}{%
    Oppgi SI-enheten og enhetssymbolet for følgende fysiske størrelser:
    \begin{enumerate}[a)]
      \item Masse
      \item Tid
      \item Temperatur
      \item Elektrisk strøm
      \item Lengde
    \end{enumerate}
  }
\end{frame}

%--------------------------------------------------------------------------
%  SLIDE 9 — Noen enheter...
%--------------------------------------------------------------------------

\begin{frame}{}
\renewcommand{\arraystretch}{1.3}
\begin{center}
\Huge{Pause!}
\end{center}
\end{frame}

%--------------------------------------------------------------------------
%  SLIDE 10 — Noen enheter...
%--------------------------------------------------------------------------

\begin{frame}{Å regne med enheter}
  \oppsum{Enhetsregning}{%
    \begin{itemize}
      \item Man bruker algebra-regler for å regne med enheter.
      \item Alle ledd i et uttrykk må ha samme enhet
    \end{itemize}
  }
  \pause
  \aktivitet{Tilbakelagt avstand}{%
En dame er på fjelltur.\\
\textbf{a)}Første etappe går hun 4 km mot nord og deretter 3 km mot vest. Hvor lang er første etappe?\\
\textbf{b)} Andre etappe holder hun 5 km/t i 1.5 timer. Hvor lang er andre etappe?
}
\end{frame}

%--------------------------------------------------------------------------
%  SLIDE 11 — Oppgave 1.2
%--------------------------------------------------------------------------

\begin{frame}{Oppgave}
  \oppgavexr{enhetsregning}{Enhetsregning}{%
    En bil kjører med fart $v = 90\,\textrm{km/t}$.
    \begin{enumerate}[a)]
      \item Gjør om farten til m/s. \textit{Hint: } $1\,\mathrm{km} = 10^{3}\,\mathrm{m}$ og $1\,\mathrm{t} = 3600\,\mathrm{s}$.
      \item Kjøreturen varer i $t = 2{,}5\,\mathrm{t}$. Finn tilbakelagt strekning i km og i m.
      \item Tyngdeakselerasjonen på månen er $g_{\mathrm{M}} = 1{,}62\,\mathrm{m/s}^{2}$. Gjør om til $\mathrm{km/min}^{2}$.
    \end{enumerate}
  }
\end{frame}

%--------------------------------------------------------------------------
%  SLIDE 12 — Behandling av tall...
%--------------------------------------------------------------------------

\begin{frame}{}
\vspace{3em}
  \centering{\Huge{Tall i fysikk!}}
\end{frame}

%--------------------------------------------------------------------------
%  SLIDE 13 — Avrunding
%--------------------------------------------------------------------------

\begin{frame}{Målenøyaktighet (signifikans) og avrunding}
\oppsum{Signifikante siffer}{%
    \begin{itemize}
        \item Antall signifikante siffer regnes fra første tall forskjellig fra null.
        \item I en utregning skal svaret angis med samme nøyaktighet (antall siffer) som den \textit{minst nøyaktige} størrelsen som inngikk i beregningen.
    \end{itemize}
}
\oppsum{Avrunding}{%
Vi minner om følgende regler for avrunding:
\begin{itemize}
    \item Tall fra 0-4 rundes ned. Tall fra og med 5 til 9 rundes opp.
    \item Avrunding til riktig antall signifikante siffer må skje helt tilslutt i en beregning!
\end{itemize}
}
\end{frame}

%--------------------------------------------------------------------------
%  SLIDE 14 — Signifikante siffer
%--------------------------------------------------------------------------
\begin{frame}{}
  \aktivitet{Beregning av areal}{%
En rektangulær hage har lengde $L = 12.3\,\mathrm{m}$ og bredde $B = 4.56\,\mathrm{m}$. Hva er arealet av hagen? Oppgi svaret med riktig antall signifikante siffer.
}


\oppgavexr{signsiff}{Signifikante siffer}{%
\begin{enumerate}[a)]
        \item Oppgi antall signifikante siffer i: $0.00420$, \ $3.14159$, \ $1\,200$ og $165.00$.
    \item Regn ut $12.3 \cdot 4.56$ og oppgi svaret med riktig antall signifikante siffer.
    \item Regn ut $18.0 / 4.5$ og oppgi svaret med riktig antall signifikante siffer.
\end{enumerate}
}
\end{frame}


\begin{frame}{Tall på standardform}
Standardform er en måte å skrive tall på som tydeliggjør både antall siffer og størrelsesorden.
\[
\textrm{T}=t \cdot 10^n,
\]
\aktivitet{Standardform}{%
\begin{enumerate}[a)]
    \item Skriv $0.000000000000000000000000001$ på standardform.
    \item Skriv $123\,000\,000\,000$ på standardform.
    \item Regn ut $1.23 \times 10^8 \div 4.56 \times 10^3$ og oppgi svaret med riktig antall signifikante siffer.
    \item Regn ut $7.89 \times 10^{-5} \cdot 1.23 \times 10^4$ og oppgi svaret med riktig antall signifikante siffer.
\end{enumerate}
}

\end{frame}

\begin{frame}{Dekadiske prefikser}
  \begin{center}
  \scriptsize
  \renewcommand{\arraystretch}{1.0}
  \begin{tabular}{|c|l|c|l|}
  \hline
  \rowcolor{blue!10!black!80}
  $10^n$ & Prefiks & Symbol & Navn \\
  \hline
  $10^{24}$ & yotta & Y & kvadrilion \\
  $10^{21}$ & zetta & Z & trilliard \\
  $10^{18}$ & exa   & E & trillion \\
  $10^{15}$ & peta  & P & billiard \\
  $10^{12}$ & tera  & T & billion \\
  $10^{9}$  & giga  & G & milliard \\
  $10^{6}$  & mega  & M & million \\
  $10^{3}$  & kilo  & k & tusen \\
  $10^{2}$  & hekto & h & hundre \\
  $10^{1}$  & deka  & da & ti \\
  $10^{-1}$ & desi  & d & tidel \\
  $10^{-2}$ & centi & c & hundredel \\
  $10^{-3}$ & milli & m & tusendel \\
  $10^{-6}$ & mikro & $\mu$ & milliondel \\
  $10^{-9}$ & nano  & n & milliarddel \\
  $10^{-12}$& piko  & p & billiondel \\
  $10^{-15}$& femto & f & billiarddel \\
  $10^{-18}$& atto  & a & trilliondel \\
  $10^{-21}$& zepto & z & trilliarddel \\
  $10^{-24}$& yocto & y & kvadrilliondel \\
  \hline
  \end{tabular}
  \end{center}
\end{frame}



%--------------------------------------------------------------------------
%  SLIDE 17 — Oppgave 1.4 og 1.5
%--------------------------------------------------------------------------

\begin{frame}{Oppgave}
  \begin{columns}[t]
    \begin{column}{0.48\textwidth}
      \oppgavexr{standardform}{Standardform}{%
        Skriv følgende tall på standardform og oppgi antall signifikante siffer:
        \begin{enumerate}[a)]
          \item $0{,}000\,045$
          \item $3\,750\,000$
          \item $0{,}100$
          \item $299\,792\,458$
        \end{enumerate}
      }
    \end{column}
    \begin{column}{0.48\textwidth}
      \oppgavexr{dekpref}{Dekadiske prefikser}{%
        \begin{enumerate}[a)]
          \item Et batteri har kapasitet $4\,200\,\mathrm{mAh}$. Skriv dette i Ah og på standardform.
          \item En avstand er $0{,}043\,\mathrm{mm}$. Skriv i $\mu\mathrm{m}$ og i nm.
          \item Lyset bruker $499\,\mathrm{s}$ fra solen. Med $c = 3{,}00\cdot10^{8}\,\mathrm{m/s}$: finn avstanden i km på standardform.
        \end{enumerate}
      }
    \end{column}
  \end{columns}
\end{frame}

\begin{frame}
\begin{center}
  \Huge{Neste gang: Kapittel 2 — Matematikk vi trenger!}
  \begin{itemize}
    \item Grunnleggende algebra
    \item Trigonometri
    \item Vektorer
  \end{itemize}
\end{center}
\end{frame}


\end{document}
